% ---------------------------------------------------
% LaTeX document generated from Markdown report
% This file is intended to be compiled with LuaLaTeX.

\documentclass[12pt]{article}

% Use modern fonts via fontspec (LuaLaTeX/XeLaTeX only)
\usepackage{fontspec}
\usepackage{geometry}
\usepackage{hyperref}
\usepackage{xcolor}
\usepackage{tikz}
\usepackage{booktabs}
\usepackage{longtable}
\usepackage{array}
\usepackage{enumitem}
\usepackage{graphicx}

% Set page geometry
\geometry{margin=1in}

% Configure hyperlink colors
\hypersetup{
  colorlinks=true,
  linkcolor=blue,
  urlcolor=blue
}

% Set a main font for LuaLaTeX; adjust as needed
\setmainfont{Latin Modern Roman}

% Macro to typeset tether IDs as footnotes safely.  It uses \detokenize
% so that the special characters inside the tether ID are not interpreted
% by LaTeX.  Each citation in the report is thus preserved as a
% footnote containing the original tether string.
\newcommand{\tcite}[1]{\footnote{\textnormal{\detokenize{#1}}}}

\begin{document}

\title{Arquitectura integral de un sistema de corrección de estilo con preservación de formato}
\author{}
\date{}
\maketitle

\tableofcontents
\newpage

\section{Introducción y objetivos}

Los libros y cartillas digitales, especialmente los que contienen diagramación, tablas e imágenes, requieren
un proceso de corrección de ortografía y de estilo que no modifique la estructura visual.  El objetivo
de este proyecto es procesar documentos largos (\textit{200--600 páginas}), corregir errores ortográficos y de
redacción de manera automatizada y mantener la fidelidad visual página por página.  Se busca evitar el uso
excesivo de servicios de terceros (APIs de OCR y LLMs en la nube) por razones de coste y privacidad.

Los documentos de entrada pueden ser \emph{PDF “born‑digital”}, \emph{PDF escaneados} o \emph{archivos DOCX}.
La salida principal es un PDF con la corrección aplicada sobre cada bloque de texto sin alterar imágenes,
tablas ni diseño.  Opcionalmente se genera un DOCX editable con correcciones.

\section{Requerimientos del sistema}

\subsection{Requerimientos funcionales}

\begin{itemize}
  \item \textbf{Ingreso de documentos}. La aplicación acepta archivos PDF y DOCX.  Si el documento
        está en formato DOCX, se convierte a PDF de forma automatizada mediante el modo sin GUI de
        \texttt{LibreOffice} (\texttt{soffice --convert-to pdf})\tcite{【333465169400446†L90-L98】}.

  \item \textbf{Extracción estructural}. Para PDFs digitales se extraen bloques de texto con
        información de posición utilizando la función \texttt{Page.get\_text("blocks")} de PyMuPDF, la
        cual devuelve una lista de bloques con \emph{coordenadas (\textit{bounding boxes})}\tcite{【915511594145204†L134-L142】}.
        Para PDFs escaneados se usa un motor de OCR que detecta y reconoce cada palabra (como
        \texttt{docTR}, que dispone de un pipeline de detección y reconocimiento de texto con prestaciones
        comparables a servicios de Google o AWS)\tcite{【191659958827309†L64-L73】}.  Para documentos muy
        complejos se usa \texttt{Layout Parser}, que permite detectar estructuras de página con modelos
        de \emph{deep learning} y abstraer las regiones en objetos \texttt{Coordinates}, \texttt{TextBlock}
        y \texttt{Layout}\tcite{【964225940706161†L47-L52】【964225940706161†L72-L76】}.

  \item \textbf{Corrección ortográfica y gramatical}. Se ejecuta un corrector local basado en
        \texttt{LanguageTool}, un corrector ortográfico, gramatical y de estilo multilingüe\tcite{【209408010763411†L104-L106】}.
        El modo ``perfeccionista'' ofrece sugerencias de estilo y tono, detecta frases largas
        o complejas y reconoce coloquialismos y redundancias\tcite{【209408010763411†L134-L139】}.

  \item \textbf{Corrección de estilo (LLM)}. Se utiliza un \emph{modelo grande de lenguaje} (LLM)
        ejecutado localmente mediante \texttt{llama.cpp} o \texttt{vLLM}.  `llama.cpp` fue diseñado
        para ejecutar modelos LLM con un montaje mínimo, con soporte de CPU + GPU y
        cuantizaciones de 1,5 a 8 bits\tcite{【949920289085484†L520-L537】}.  Los modelos permiten generar
        parches de estilo (pequeñas sustituciones) para cada bloque sin reescribir todo el contenido.

  \item \textbf{Control de desbordamiento}.  Cada corrección se vuelve a insertar en su
        \emph{bounding box} y se valida que el texto corregido quepa.  Si no cabe, se aplican
        reglas de ajuste (reducción de tamaño de fuente o espaciado) o se marca para revisión
        manual.

  \item \textbf{Continuidad y coherencia}. El sistema crea \emph{instantáneas de contexto} que almacenan
        glosarios, preferencias de estilo y entidades detectadas hasta la página \(N\) para garantizar
        consistencia en terminología y tono en páginas siguientes.

  \item \textbf{Interfaz de usuario}. La aplicación ofrece una vista de diferencias lado a lado
        (original vs corregido) con resaltado por tipo de cambio y botones para aceptar o
        rechazar correcciones.  Permite ajustar reglas de estilo, diccionarios personalizados y
        glosarios por proyecto.

  \item \textbf{Salida final}. Se genera un PDF con el texto corregido sobreimpreso conservando
        imágenes y diagramación.  Opcionalmente se genera un DOCX editable para futuras ediciones.
\end{itemize}

\subsection{Requerimientos no funcionales}

\begin{itemize}
  \item \textbf{Precisión y fidelidad}.  El texto debe corregirse sin introducir nuevos errores ni
        alterar el significado.  La corrección no debe desordenar el diseño de página.

  \item \textbf{Escalabilidad}.  El sistema debe procesar documentos de cientos de páginas en
        paralelo.  Se usan tareas asincrónicas mediante \texttt{Celery}; Celery es un sistema
        distribuido que procesa grandes cantidades de mensajes, con enfoque en procesamiento en
        tiempo real y soporte para la programación de tareas\tcite{【227726339796927†L13-L19】}.

  \item \textbf{Privacidad y coste}.  La ejecución de OCR y LLM debe realizarse en servidores
        controlados por la organización para reducir costes y proteger datos sensibles.

  \item \textbf{Extensibilidad}.  El sistema debe permitir la sustitución de componentes (por
        ejemplo, cambiar el motor de OCR o LLM) y la personalización de reglas de estilo sin
        reescribir todo el pipeline.

  \item \textbf{Tolerancia a fallos}.  Debe permitir recuperar procesos en caso de fallo de
        hardware o software.  Los artefactos intermedios se almacenan en un repositorio de objetos
        y la base de datos mantiene el estado de cada etapa.

  \item \textbf{Seguridad}.  Los archivos y metadatos deben almacenarse de forma cifrada; el
        acceso a la interfaz debe estar protegido mediante autenticación y control de roles.
\end{itemize}

\section{Arquitectura general}

\subsection{Visión de alto nivel}

La aplicación se construye siguiendo una \emph{arquitectura de microservicios} con un front‑end y un
back‑end desacoplados.  El sistema se compone de los siguientes bloques principales:

\begin{enumerate}[label=\arabic*.]
  \item \textbf{UI Web} (Next.js / React) – Permite subir documentos, mostrar el progreso,
        revisar correcciones y descargar resultados.
  \item \textbf{API} (FastAPI) – Maneja peticiones del cliente, autentica usuarios y orquesta
        tareas.
  \item \textbf{Cola de tareas} (Celery con Redis o RabbitMQ) – Gestiona trabajos pesados
        (extracción, OCR, corrección, renderizado).  Celery ofrece un sistema de colas
        distribuido y flexible\tcite{【227726339796927†L13-L19】}.
  \item \textbf{Workers} – Conjunto de servicios que ejecutan etapas de procesamiento (ingesta,
        extracción, corrección, renderizado) de forma paralela.
  \item \textbf{Almacenamiento de objetos} (MinIO) – Se usa como repositorio S3 compatible para
        almacenar documentos originales, artefactos intermedios, salidas y datos binarios.
        MinIO es un almacenamiento de objetos de alto rendimiento compatible con S3, diseñado
        para velocidad y escalabilidad en cargas de trabajo de IA/ML y análisis\tcite{【870299264509201†L414-L421】}.
  \item \textbf{Base de datos} (PostgreSQL) – Gestiona el modelo de datos relacional y el estado
        de cada documento.  PostgreSQL es un sistema de base de datos objeto‑relacional de código
        abierto con más de 35 años de desarrollo y reputación de fiabilidad, robustez y
        rendimiento\tcite{【357047729773606†L23-L25】}.
  \item \textbf{Módulos de procesamiento}:
    \begin{itemize}
      \item \emph{Conversión a PDF} (LibreOffice headless) – Convierte DOCX a PDF sin interfaz
            gráfica con el comando \texttt{soffice --convert-to pdf}, que activa el modo
            headless\tcite{【333465169400446†L90-L98】}.
      \item \emph{Extracción de layout} – Usa PyMuPDF para obtener bloques de texto y sus
            posiciones en PDFs digitales.  En PDFs escaneados se usa docTR para OCR de dos
            etapas\tcite{【191659958827309†L64-L73】} y Layout Parser para detectar estructuras de página
            mediante modelos de \emph{deep learning}\tcite{【964225940706161†L47-L52】}.
      \item \emph{Corrección determinista} – \texttt{LanguageTool} local para ortografía, gramática y
            estilo\tcite{【209408010763411†L104-L106】}.
      \item \emph{Corrección LLM} – Un LLM local (por ejemplo, Qwen/Mistral) ejecutado con
            \texttt{llama.cpp}.  `llama.cpp` permite inferencia de modelos grandes en C/C++ con
            soporte de CPU y GPU, con optimizaciones por plataforma y cuantización de 1,5 a 8 bits\tcite{【949920289085484†L520-L537】}.
      \item \emph{Renderizado} – Reescribe el texto en su misma caja mediante PyMuPDF,
            validando la cabida y ajustando la fuente si es necesario.
    \end{itemize}
\end{enumerate}

\subsection{Diagrama de componentes}

\begin{figure}[h!]
  \centering
  \begin{tikzpicture}[node distance=1.5cm, every node/.style={font=\small}]
    % Define component nodes
    \node (client) [draw, rectangle, rounded corners, fill=blue!10, minimum width=3cm, minimum height=0.8cm] {Cliente Web};
    \node (api) [draw, rectangle, rounded corners, fill=green!10, below=of client] {API FastAPI};
    \node (db) [draw, rectangle, rounded corners, fill=yellow!10, right=3cm of api] {PostgreSQL};
    \node (queue) [draw, rectangle, rounded corners, fill=orange!10, below=of api] {Cola Celery};
    \node (workers) [draw, rectangle, rounded corners, fill=red!10, below=of queue] {Workers};
    \node (storage) [draw, rectangle, rounded corners, fill=purple!10, right=3cm of workers] {Almacenamiento MinIO};
    
    % Draw arrows between components
    \draw[->] (client) -- (api);
    \draw[->] (api) -- (queue);
    \draw[->] (api) -- (db);
    \draw[->] (queue) -- (workers);
    \draw[->] (workers) -- (storage);
    \draw[->] (workers) |- (db);

    % Legend or labels if needed
    \node[above=0.3cm of client] {\textbf{Front‑end}};
    \node[above=0.3cm of api] {\textbf{Back‑end}};
  \end{tikzpicture}
  \caption{Diagrama de componentes de alto nivel.}
  \label{fig:componentes}
\end{figure}

En la Figura \ref{fig:componentes} se representa la interacción entre los componentes principales del sistema. El
usuario carga un archivo a través de la interfaz web; la API gestiona las peticiones y guarda el
archivo en MinIO, además de registrar la entrada en la base de datos.  Los trabajos pesados se envían a
la cola de Celery, que los distribuye entre los trabajadores.  Estos procesan cada etapa y almacenan
los artefactos en MinIO y el estado en Postgres.

\section{Flujo de procesamiento}

El sistema se organiza en etapas encadenadas para asegurar la fidelidad y la coherencia.  Cada etapa
produce artefactos intermedios que se almacenan y se asocian a un registro en la base de datos.

\subsection{Etapa A: Ingesta}

\begin{enumerate}[label=\arabic*.]
  \item \textbf{Subida y normalización}. El documento original se guarda en \texttt{source/}
        en MinIO y se registra en la tabla \texttt{documents}.
  \item \textbf{Conversión a PDF}.  Si el archivo es DOCX u otro formato soportado se
        convierte a PDF con LibreOffice headless (\texttt{soffice --convert-to pdf})\tcite{【333465169400446†L90-L98】}.
        El PDF convertido se almacena en \texttt{pdf/}.
  \item \textbf{Conteo de páginas}. Se calcula el número de páginas y se crean registros
        \texttt{pages} con estado inicial.
\end{enumerate}

\subsection{Etapa B: Extracción de layout}

Se ejecuta un \emph{worker} por página:

\begin{itemize}
  \item Para PDFs digitales, \textbf{PyMuPDF} analiza cada página y extrae bloques de texto
        (\texttt{get\_text("blocks")}) con coordenadas, fuentes y estilos\tcite{【915511594145204†L134-L142】}.  También se
        extraen imágenes y tablas.  Cada página genera un JSON \texttt{layout} con campos
        \{\texttt{bbox}, tipo, spans\} y una imagen de vista previa.

  \item Para PDFs escaneados, \textbf{docTR} realiza detección y reconocimiento de texto en
        dos etapas\tcite{【191659958827309†L64-L73】}.  Se obtienen cajas de texto (palabras) y su
        contenido.

  \item Cuando las páginas contienen estructuras complejas, \textbf{Layout Parser} ejecuta
        modelos de \emph{deep learning} para detectar regiones (cabeceras, párrafos, tablas,
        figuras) en pocas líneas de código y sin reglas fijas\tcite{【964225940706161†L47-L52】}.  El
        resultado se almacena en la abstracción \texttt{Layout}, compuesta por objetos
        \texttt{Coordinates}, \texttt{TextBlock} y \texttt{Layout}\tcite{【964225940706161†L72-L76】}.

  \item Cada página almacena sus recursos (layout JSON, texto plano, imágenes) en MinIO.
\end{itemize}

\subsection{Etapa C: Construcción de contexto}

Una vez procesada una ventana de \(N\) páginas (por ejemplo, 10 páginas), se crea una instantánea de
contexto que almacena glosarios, entidades, preferencias de estilo y resumen del hilo para
mantener coherencia a lo largo del documento.  Este contexto se actualiza de manera incremental:

\begin{itemize}
  \item Se combinan las entidades y términos detectados en las páginas previas.
  \item Se aplican reglas de estilo (tuteo/usted, marcas de cita) y glosarios personalizados.
  \item Se almacena un resumen de la narrativa para que las correcciones posteriores respeten
        continuidad temática.
\end{itemize}

El contexto se asocia a un campo \texttt{snapshot} en la base de datos y cada corrección lo utiliza
como referencia.

\subsection{Etapa D: Corrección por página}

El proceso de corrección se aplica por bloque dentro de cada página.  Se controla el desbordamiento
para que el texto corregido quepa en su caja original.  Los pasos son los siguientes:

\begin{enumerate}[label=\arabic*.]
  \item \textbf{Pre‑chequeo determinista}.  Se ejecuta \texttt{LanguageTool} sobre cada bloque
        de texto para corregir ortografía, gramática y puntuación.
  \item \textbf{Corrección de estilo con LLM}.  Un modelo LLM local revisa el bloque y genera
        parches de estilo (reemplazos, inserciones, eliminaciones) acotados.  Las respuestas
        se devuelven como una estructura JSON de operaciones y no como texto libre.
  \item \textbf{Validación de cabida}.  Se inserta el texto corregido en su bounding box usando
        PyMuPDF.  Si el nuevo texto no cabe:
    \begin{enumerate}[label=\alph*.]
      \item Se aplica una reducción de tamaño de fuente o un ajuste de espaciado hasta un
            umbral (por ejemplo, −10 %).
      \item Si la reducción no es suficiente, se marca el bloque para revisión manual.
    \end{enumerate}
  \item \textbf{Post‑chequeo}.  Se vuelve a ejecutar \texttt{LanguageTool} para asegurar que no
        queden errores.  Se calcula una métrica de calidad (\textit{qa\_score}) y se registra el
        \emph{flag} de desbordamiento.
  \item \textbf{Guardar parches}.  El parche y las métricas se guardan en la tabla
        \texttt{patches} y en la ruta correspondiente en MinIO.
\end{enumerate}

\subsection{Etapa E: Renderizado y ensamble}

Finalmente se generan las páginas del PDF final.  Para cada página se aplica el parche
validado sobre la capa de texto original usando PyMuPDF; se reinsertan imágenes y tablas
sin modificación.  Una vez procesadas todas las páginas, se ensamblan en un único PDF.

El flujo completo se resume en el diagrama de la Figura~\ref{fig:pipeline}, donde cada
etapa se representa como un bloque en una cadena y se muestra la interacción con el
almacenamiento y la base de datos.

\begin{figure}[h!]
  \centering
  \begin{tikzpicture}[node distance=2.2cm, every node/.style={font=\small}, align=center]
    % Define nodes for pipeline stages
    \node (ingest) [draw, rectangle, rounded corners, fill=blue!10, minimum width=3.2cm] {Etapa A\\Ingesta};
    \node (extract) [draw, rectangle, rounded corners, fill=green!10, right=of ingest] {Etapa B\\Extracción de layout};
    \node (context) [draw, rectangle, rounded corners, fill=yellow!10, right=of extract] {Etapa C\\Construcción de contexto};
    \node (correct) [draw, rectangle, rounded corners, fill=orange!10, right=of context] {Etapa D\\Corrección};
    \node (render) [draw, rectangle, rounded corners, fill=red!10, right=of correct] {Etapa E\\Renderizado};

    % Arrows connecting pipeline stages
    \draw[->] (ingest) -- (extract);
    \draw[->] (extract) -- (context);
    \draw[->] (context) -- (correct);
    \draw[->] (correct) -- (render);

    % Storage and DB on the side
    \node (storage2) [draw, rectangle, rounded corners, fill=purple!10, below=1.2cm of extract, minimum width=3cm] {MinIO};
    \node (db2) [draw, rectangle, rounded corners, fill=pink!10, below=1.2cm of correct, minimum width=3cm] {PostgreSQL};

    % Arrows to storage and DB
    \draw[->] (ingest.south) -- (storage2.north west);
    \draw[->] (extract.south) -- (storage2.north);
    \draw[->] (context.south) -- (db2.north west);
    \draw[->] (correct.south) -- (db2.north);
    \draw[->] (render.south) -- (storage2.north east);

    % Label arrows to indicate artifacts
    \node at ($(ingest.south)!0.5!(storage2.north west)$) [left=0.1cm] {\scriptsize artefactos};
    \node at ($(extract.south)!0.5!(storage2.north)$) {\scriptsize layouts/imagenes};
    \node at ($(context.south)!0.5!(db2.north west)$) {\scriptsize snapshots};
    \node at ($(correct.south)!0.5!(db2.north)$) {\scriptsize parches};
    \node at ($(render.south)!0.5!(storage2.north east)$) {\scriptsize PDF final};
  \end{tikzpicture}
  \caption{Flujo de procesamiento y generación de artefactos.}
  \label{fig:pipeline}
\end{figure}

\section{Modelo de datos y almacenamiento}

El sistema utiliza una base de datos relacional (PostgreSQL) junto con un repositorio de objetos
(MinIO) para gestionar tanto los metadatos como los contenidos binarios.  El modelo de datos se
resume en la Tabla~\ref{tab:datos}, donde cada tabla principal se describe con sus campos más
relevantes.

\begin{longtable}{>{\raggedright\arraybackslash}p{3cm}p{4.5cm}p{6.5cm}}
  \caption{Tablas principales del modelo de datos.}\label{tab:datos}\\
  \toprule
  \textbf{Tabla} & \textbf{Campos (resumen)} & \textbf{Descripción} \\
  \midrule
  \endfirsthead
  \caption[]{Tablas principales del modelo de datos (continuación).}\\
  \toprule
  \textbf{Tabla} & \textbf{Campos (resumen)} & \textbf{Descripción} \\
  \midrule
  \endhead
  \bottomrule
  \endfoot
  \texttt{documents} & \texttt{id}, \texttt{filename}, \texttt{source\_uri}, \texttt{pdf\_uri}, \texttt{total\_pages}, \texttt{status}, \texttt{config\_json}, \texttt{created\_at} & Registro de cada documento. \texttt{config\_json} almacena reglas de estilo, idioma y preferencias. \\
  \addlinespace
  \texttt{pages} & \texttt{id}, \texttt{doc\_id}, \texttt{page\_no}, \texttt{layout\_uri}, \texttt{text\_uri}, \texttt{status}, \texttt{preview\_uri} & Cada página del documento. Guarda su layout JSON, el texto plano y la miniatura. \\
  \addlinespace
  \texttt{patches} & \texttt{id}, \texttt{page\_id}, \texttt{version}, \texttt{patch\_uri}, \texttt{qa\_score}, \texttt{overflow\_flag}, \texttt{applied} & Cambios propuestos por LLM y LanguageTool. \texttt{applied} indica si se usó para renderizar; \texttt{qa\_score} mide la calidad. \\
  \addlinespace
  \texttt{context\_snapshots} & \texttt{id}, \texttt{doc\_id}, \texttt{upto\_page\_no}, \texttt{snapshot\_uri}, \texttt{created\_at} & Instantáneas de contexto acumulado hasta cierta página.  Incluye glosarios, entidades y preferencias. \\
  \addlinespace
  \texttt{jobs} & \texttt{id}, \texttt{doc\_id}, \texttt{task\_type}, \texttt{status}, \texttt{started\_at}, \texttt{finished\_at}, \texttt{error} & Registros de trabajos en Celery: ingesta, extracción, corrección y renderizado. \\
\end{longtable}

\subsection{Almacenamiento en objetos (MinIO)}

El contenido binario se organiza en una estructura de carpetas jerárquica en MinIO.  Cada
documento tiene su propio espacio y las rutas se dividen de la siguiente manera:

\begin{itemize}
  \item \texttt{source/} – Documentos originales cargados.
  \item \texttt{pdf/} – Documentos convertidos a PDF (en caso de entrada DOCX).
  \item \texttt{pages/layout/} – Archivos JSON con la estructura de cada página.
  \item \texttt{pages/text/} – Texto plano de cada página.
  \item \texttt{pages/preview/} – Imágenes de vista previa para la interfaz.
  \item \texttt{pages/patch/} – Parches generados por la IA.
  \item \texttt{pages/output/} – Fragmentos de páginas después de la renderización.
  \item \texttt{final/} – Documentos finales listos para descarga.
\end{itemize}

MinIO es un almacenamiento de objetos de alto rendimiento, compatible con la API S3 y optimizado para
cargas de trabajo de IA y análisis\tcite{【870299264509201†L414-L421】}.  Por tanto, es adecuado para manejar
grandes volúmenes de archivos y escala horizontalmente.

\section{Tecnologías y justificación}

En la Tabla~\ref{tab:tecnologias} se listan los componentes principales, la tecnología
seleccionada y la justificación basada en las fuentes consultadas.

\begin{longtable}{>{\raggedright\arraybackslash}p{3.5cm}p{3.5cm}p{6cm}}
  \caption{Tecnologías utilizadas y justificación.}\label{tab:tecnologias}\\
  \toprule
  \textbf{Componente} & \textbf{Tecnología} & \textbf{Justificación y cita} \\
  \midrule
  \endfirsthead
  \caption[]{Tecnologías utilizadas y justificación (continuación).}\\
  \toprule
  \textbf{Componente} & \textbf{Tecnología} & \textbf{Justificación y cita} \\
  \midrule
  \endhead
  \bottomrule
  \endfoot
  Extracción de PDFs & \textbf{PyMuPDF} & Permite obtener bloques de texto y sus bounding boxes mediante \texttt{Page.get\_text("blocks")}\tcite{【915511594145204†L134-L142】} y extrae imágenes y tablas. \\
  \addlinespace
  OCR de escaneados & \textbf{docTR} (Mindee) & Librería con pipeline de detección y reconocimiento de texto en dos etapas, con rendimiento comparable a servicios comerciales\tcite{【191659958827309†L64-L73】}. \\
  \addlinespace
  Detección de layout & \textbf{Layout Parser} & Toolkit de \emph{deep learning} para extraer estructuras complejas de documentos con pocas líneas de código y sin reglas fijas\tcite{【964225940706161†L47-L52】}; ofrece estructuras \texttt{Coordinates}, \texttt{TextBlock} y \texttt{Layout}\tcite{【964225940706161†L72-L76】}. \\
  \addlinespace
  Conversión de DOCX & \textbf{LibreOffice headless} & Permite convertir DOCX a PDF sin interfaz gráfica, activando automáticamente el modo headless cuando se usa \texttt{--convert-to pdf}\tcite{【333465169400446†L90-L98】}. \\
  \addlinespace
  Corrección determinista & \textbf{LanguageTool local} & Corrector ortográfico, gramatical y de estilo multilingüe\tcite{【209408010763411†L104-L106】} que sugiere mejoras de estilo y detecta frases largas\tcite{【209408010763411†L134-L139】}. \\
  \addlinespace
  Corrección LLM & \textbf{llama.cpp / vLLM} & `llama.cpp` permite ejecutar modelos LLM localmente con mínima configuración y soporte para CPU+GPU, con optimizaciones y cuantización de 1,5 hasta 8 bits\tcite{【949920289085484†L520-L537】}. \\
  \addlinespace
  Orquestación de tareas & \textbf{Celery} & Cola de tareas distribuida enfocada en procesamiento en tiempo real y programación de tareas\tcite{【227726339796927†L13-L19】}. \\
  \addlinespace
  API backend & \textbf{FastAPI} (Python) & Framework asíncrono para construir APIs rápidas, con tipado, validación y
        integración sencilla con Pydantic. \\
  \addlinespace
  Frontend & \textbf{Next.js / React} & Construye una interfaz reactiva con vista de diferencias, permite SSR y generación estática. \\
  \addlinespace
  Base de datos & \textbf{PostgreSQL} & Sistema de base de datos relacional robusto y de alto rendimiento\tcite{【357047729773606†L23-L25】}. \\
  \addlinespace
  Almacenamiento & \textbf{MinIO} & Almacenamiento de objetos de alto rendimiento, compatible con S3 y optimizado para IA/ML\tcite{【870299264509201†L414-L421】}. \\
  \addlinespace
  Contenedor y despliegue & \textbf{Docker + Kubernetes} & Permiten empaquetar servicios y escalar workers de manera horizontal; se pueden aprovechar GPUs para OCR y LLM. \\
  \addlinespace
  Mensajería interna & \textbf{Redis} o RabbitMQ & Servidores de cola utilizados por Celery para distribuir trabajos a los workers. \\
  \addlinespace
  Caché & \textbf{Redis} & Cache para contextos y resultados de consultas frecuentes. \\
  \addlinespace
  Monitorización & \textbf{Prometheus + Grafana} & Recopilan métricas de workers, uso de memoria y tiempos de cada etapa; permiten ajustar el dimensionamiento. \\
\end{longtable}

\section{Despliegue y escalado}

Para garantizar un funcionamiento robusto y un escalado flexible, el despliegue utiliza contenedores
orquestados por Kubernetes.  Las consideraciones principales son:

\begin{itemize}
  \item \textbf{Contenedores}.  Cada componente (API, workers, Celery, MinIO, Redis, PostgreSQL)
        se despliega en contenedores individuales.
  \item \textbf{Kubernetes}.  Se utiliza un orquestador para escalar horizontalmente los workers
        según la carga y administrar reinicios automáticos.
  \item \textbf{GPU opcional}.  Los workers de OCR y LLM pueden tener acceso a GPU; Kubernetes
        programa pods en nodos con GPU mediante etiquetas.
  \item \textbf{Persistencia}.  MinIO se ejecuta en modo distribuido (al menos 4 nodos para
        paridad) para tolerancia a fallos.  PostgreSQL se despliega con réplicas y almacenamiento
        persistente.
  \item \textbf{Seguridad}.  Se definen roles y políticas de acceso en API y MinIO (por ejemplo,
        bucket‑policies S3).  Se utiliza TLS para todas las comunicaciones y se almacena el
        modelo de credenciales en un gestor de secretos (Vault/Secrets).
  \item \textbf{Actualizaciones}.  Los modelos de IA y los componentes se versionan y se
        despliegan con \emph{rolling updates} para minimizar tiempos de inactividad.
\end{itemize}

\section{Consideraciones de rendimiento y calidad}

\begin{itemize}
  \item \textbf{Optimización de OCR}.  Para documentos mixtos, se utiliza OCR solo en
        páginas escaneadas.  \texttt{docTR} es eficiente y optimizado tanto para CPU como GPU\tcite{【191659958827309†L64-L73】}.

  \item \textbf{Paralelismo por ventana}.  Las páginas se dividen en ventanas (por ejemplo, 10 páginas)
        para procesar varias en paralelo mientras se mantiene el contexto entre ventanas.

  \item \textbf{Cuantización de LLM}.  La cuantización a 4 u 8 bits de los modelos LLM reduce la
        memoria y acelera la inferencia; `llama.cpp` soporta cuantizaciones de 1,5 hasta 8 bits\tcite{【949920289085484†L520-L537】}.

  \item \textbf{Control de desbordamiento}.  Se establecen límites de reducción de fuente (por ejemplo,
        hasta –10 %) para garantizar que el texto siga siendo legible y que los cambios no
        distorsionen la paginación.

  \item \textbf{Métricas}.  Se definen métricas como: número de errores detectados y corregidos,
        tasa de desbordamiento, tiempo de procesamiento por página y número de revisiones manuales
        necesarias.
\end{itemize}

\section{Seguridad y privacidad}

\begin{itemize}
  \item \textbf{Aislamiento de datos}.  Los documentos se almacenan en MinIO con cifrado en
        reposo y se separan por cliente.

  \item \textbf{Cifrado en tránsito}.  Todas las comunicaciones (API, Celery, MinIO) usan TLS.

  \item \textbf{Control de acceso}.  Los usuarios tienen roles (editor, revisor, administrador).
        Los tokens o claves se rotan periódicamente.

  \item \textbf{Auditoría}.  Se registra cada acción en la base de datos y se pueden generar
        trazas por bloque (qué se cambió y por qué).

  \item \textbf{Cumplimiento}.  Se pueden activar configuraciones de MinIO para inmovilizar
        objetos e impedir su borrado durante períodos definidos (retención)\tcite{【870299264509201†L414-L421】}.
\end{itemize}

\section{Mantenimiento y extensibilidad}

\begin{itemize}
  \item \textbf{Actualización de modelos}.  Los modelos de OCR y LLM se actualizan de forma
        modular.  Los datos de entrenamiento y los glosarios se almacenan para realizar
        \emph{fine‑tuning} si se dispone de corpus corregidos.

  \item \textbf{Personalización de reglas}.  El sistema permite añadir reglas de estilo
        específicas (por ejemplo, usar siempre “documento” en lugar de “dossier”) y glosarios
        por proyecto.

  \item \textbf{Internacionalización}.  \texttt{LanguageTool} admite más de 25 idiomas\tcite{【209408010763411†L104-L119】}; se pueden instalar recursos adicionales.

  \item \textbf{API externa}.  La aplicación expone endpoints REST y GraphQL para integrarse
        con sistemas externos.

  \item \textbf{Documentación y pruebas}.  El backend se documenta con OpenAPI/Swagger.
        Se incluyen pruebas unitarias y de integración para cada etapa del pipeline.
\end{itemize}

\section{Conclusión}

Este documento presenta una \emph{arquitectura integral para un sistema de corrección de estilo} que
preserva la estructura visual de documentos complejos.  Combina tecnologías de extracción
(PyMuPDF, Layout Parser, docTR) con correctores deterministas (LanguageTool) y LLMs locales
ejecutados mediante `llama.cpp`.  El diseño modular y la orquestación asíncrona permiten procesar
libros de cientos de páginas con fidelidad, garantizando escalabilidad y control de costes.

La adopción de almacenamiento de objetos S3‑compatible (MinIO)\tcite{【870299264509201†L414-L421】}, una base de datos robusta
(PostgreSQL)\tcite{【357047729773606†L23-L25】} y una cola distribuida (Celery)\tcite{【227726339796927†L13-L19】} asegura que el sistema
sea extensible, seguro y sostenible.

\end{document}